\chapter{Funkcje ciągłe}
\label{ma1:ch:continuous-functions}

Granica odwzorowania w punkcie skupienia dziedziny opisuje zachowanie jego wartości w przekłutym otoczeniu. Ciągłość wymaga zgodności tego zachowania z wartością odwzorowania w samym punkcie. Definicja ciągłości ma jednak sens także w punktach izolowanych. W tym rozdziale wszystkie podzbiory przestrzeni metrycznych rozpatrujemy z metryką odziedziczoną, o ile nie wskazano inaczej.

\section{Definicja i kryteria ciągłości}
\begin{definicja}
\label{ma1:def:continuity}
Niech $(X,d_X)$ i $(Y,d_Y)$ będą przestrzeniami metrycznymi, $A\subseteq X$, $f:A\to Y$ i $p\in A$. Odwzorowanie $f$ jest ciągłe w punkcie $p$, jeżeli
\[
 \forall\varepsilon>0\ \exists\delta>0\ \forall x\in A:
 \quad d_X(x,p)<\delta\ \Longrightarrow\ d_Y(f(x),f(p))<\varepsilon.
\]
Odwzorowanie $f$ jest ciągłe na $A$, jeżeli jest ciągłe w każdym punkcie $A$.
\end{definicja}
W odróżnieniu od definicji granicy nie wykluczamy $x=p$; wówczas odległość obrazów wynosi zero. Jeśli $p\in A\cap A'$, ciągłość w $p$ jest równoważna równości $\lim_{A\ni x\to p}f(x)=f(p)$. Jeżeli $p\in A\setminus A'$, istnieje $\delta>0$ takie, że jedynym punktem $A$ w kuli o środku $p$ i promieniu $\delta$ jest $p$. Każde odwzorowanie $f:A\to Y$ jest wtedy ciągłe w $p$.

\begin{przyklad}
\label{ma1:ex:isolated-continuity}
Każde odwzorowanie $f:\Z\to\R$ jest ciągłe: dla $p\in\Z$ i $|x-p|<1/2$, $x\in\Z$, mamy $x=p$. Nie jest potrzebne istnienie granicy w przekłutym otoczeniu.
Natomiast funkcja z przykładu \ref{ma1:ex:removable-value} nie jest ciągła w zerze, ponieważ jej granica wynosi $1$, a wartość $7$. Zmiana wyłącznie wartości w zerze na $1$ daje funkcję ciągłą na $\R$.
\end{przyklad}

\begin{twierdzenie}[Kryterium ciągowe ciągłości]
\label{ma1:thm:continuity-sequential}
Niech $X,Y$ będą przestrzeniami metrycznymi, $A\subseteq X$, $p\in A$ i $f:A\to Y$. Odwzorowanie $f$ jest ciągłe w $p$ wtedy i tylko wtedy, gdy dla każdego ciągu $(x_n)$ w $A$ zbieżnego do $p$ zachodzi $f(x_n)\to f(p)$.
\end{twierdzenie}
\begin{proof}
Jeżeli $f$ jest ciągłe w $p$, to dla $\varepsilon>0$ wybieramy $\delta>0$ z definicji. Z $x_n\to p$ wynika, że od pewnego indeksu $d_X(x_n,p)<\delta$, stąd $d_Y(f(x_n),f(p))<\varepsilon$.

Odwrotnie, zaprzeczenie ciągłości oznacza istnienie $\varepsilon_0>0$ takiego, że dla każdego $\delta>0$ istnieje $x\in A$ spełniający $d_X(x,p)<\delta$ oraz $d_Y(f(x),f(p))\ge\varepsilon_0$. Dla $\delta=1/n$, $n\ge1$, wybieramy taki $x_n$. Wówczas $x_n\to p$, lecz $f(x_n)$ nie dąży do $f(p)$, sprzecznie z założeniem.
\end{proof}

\section{Zbiory otwarte i domknięte}
\begin{definicja}
\label{ma1:def:relative-open}
Niech $(X,d)$ będzie przestrzenią metryczną i $A\subseteq X$. Zbiór $U\subseteq A$ nazywamy otwartym w $A$, jeżeli dla każdego $p\in U$ istnieje $r>0$ takie, że $B_X(p,r)\cap A\subseteq U$, gdzie $B_X(p,r)=\{x\in X:d(x,p)<r\}$. Zbiór $F\subseteq A$ nazywamy domkniętym w $A$, jeżeli $A\setminus F$ jest otwarty w $A$.
\end{definicja}
Dla $A=X$ otrzymujemy zwykłą definicję zbioru otwartego i równoważną definicję zbioru domkniętego z rozdziału \ref{ma1:ch:cauchy-completeness}: jeśli każdy ciąg w $F$ zbieżny w $X$ ma granicę w $F$, to dla $p\notin F$ pewna kula nie przecina $F$, gdyż inaczej moglibyśmy wybrać $x_n\in F\cap B_X(p,1/n)$, sprzecznie z własnością granic. Odwrotnie, gdy $X\setminus F$ jest otwarty, ciąg w $F$ nie może dążyć do punktu tego dopełnienia, ponieważ od pewnego indeksu należałby do kuli rozłącznej z $F$. To samo rozumowanie stosuje się w przestrzeni $A$.

\begin{lemat}
\label{ma1:lem:relative-open-representation}
Niech $(X,d)$ będzie przestrzenią metryczną, $A\subseteq X$ oraz $U,F\subseteq A$. Zbiór $U$ jest otwarty w $A$ wtedy i tylko wtedy, gdy istnieje zbiór otwarty $V\subseteq X$ taki, że $U=A\cap V$. Analogicznie $F$ jest domknięty w $A$ wtedy i tylko wtedy, gdy istnieje zbiór domknięty $H\subseteq X$ taki, że $F=A\cap H$.
\end{lemat}
\begin{proof}
Jeżeli $U$ jest otwarty w $A$, to dla każdego $p\in U$ istnieje $r_p>0$ takie, że $B_X(p,r_p)\cap A\subseteq U$. Niech $V$ będzie sumą wszystkich takich kul (dla $U=\varnothing$ przyjmujemy $V=\varnothing$). Każda kula jest otwarta: jeśli $y\in B_X(p,r_p)$, to kula $B_X(y,r_p-d(y,p))$ zawiera się w $B_X(p,r_p)$ na mocy nierówności trójkąta. Dowolna suma zbiorów otwartych jest otwarta, ponieważ każdy punkt sumy należy do pewnego jej składnika i ma w nim otoczenie. Zatem $V$ jest otwarty w $X$. Każdy $p\in U$ należy do $A\cap V$, a każda kula użyta w konstrukcji $V$ przecina $A$ wyłącznie w punktach $U$. Stąd $U=A\cap V$.

Odwrotnie, jeżeli $U=A\cap V$, gdzie $V$ jest otwarty w $X$, to dla każdego $p\in U$ istnieje $r>0$ takie, że $B_X(p,r)\subseteq V$. W konsekwencji $B_X(p,r)\cap A\subseteq U$. Dla zbiorów domkniętych stosujemy udowodnioną równoważność do $A\setminus F$ oraz tożsamość $A\setminus(A\cap V)=A\cap(X\setminus V)$.
\end{proof}

\begin{przyklad}
\label{ma1:ex:relative-open}
Zbiór $[0,1/2)$ jest otwarty w $A=[0,1]$, gdyż jest równy $A\cap(-1,1/2)$. Nie jest otwarty w $\R$, ponieważ każda kula o środku $0$ zawiera liczby ujemne. Zbiór $A$ jest zarazem otwarty i domknięty w samym sobie, podobnie jak zbiór pusty.
\end{przyklad}

\begin{twierdzenie}
\label{ma1:thm:continuity-preimages}
Niech $X,Y$ będą przestrzeniami metrycznymi, $A\subseteq X$, $f:A\to Y$. Następujące warunki są równoważne:
\begin{enumerate}[label=(\roman*)]
\item $f$ jest ciągłe na $A$;
\item dla każdego otwartego zbioru $V\subseteq Y$ przeciwobraz $f^{-1}(V)=\{x\in A:f(x)\in V\}$ jest otwarty w $A$;
\item dla każdego domkniętego zbioru $F\subseteq Y$ przeciwobraz $f^{-1}(F)$ jest domknięty w $A$.
\end{enumerate}
\end{twierdzenie}
\begin{proof}
(i)$\Rightarrow$(ii). Niech $p\in f^{-1}(V)$. Ponieważ $V$ jest otwarty i $f(p)\in V$, istnieje $\varepsilon>0$ takie, że $B_Y(f(p),\varepsilon)\subseteq V$. Ciągłość daje $\delta>0$ takie, że $f(B_X(p,\delta)\cap A)\subseteq B_Y(f(p),\varepsilon)$. Stąd $B_X(p,\delta)\cap A\subseteq f^{-1}(V)$, co dowodzi otwartości.

(ii)$\Rightarrow$(i). Dla $p\in A$ oraz $\varepsilon>0$ kula $V=B_Y(f(p),\varepsilon)$ jest otwarta: dla $y\in V$ kula o środku $y$ i promieniu $\varepsilon-d_Y(y,f(p))>0$ zawiera się w $V$ na mocy nierówności trójkąta. Z założenia $f^{-1}(V)$ jest otwarty w $A$ i zawiera $p$. Istnieje więc $\delta>0$ z $B_X(p,\delta)\cap A\subseteq f^{-1}(V)$. Jest to warunek ciągłości.

(ii)$\Leftrightarrow$(iii). Dla każdego $V\subseteq Y$ mamy $A\setminus f^{-1}(V)=f^{-1}(Y\setminus V)$, gdyż obie strony składają się dokładnie z punktów, których obrazy nie należą do $V$. Stosujemy definicję zbioru domkniętego w obu przestrzeniach.
\end{proof}
Nie należy zastępować przeciwobrazów obrazami: ciągłe odwzorowanie stałe $\R\to\R$, $x\mapsto0$, przekształca otwarty przedział w nieotwarty singleton. Z kolei ciągłe odwzorowanie $x\mapsto 1/(1+x^2)$ przekształca domknięty zbiór $\R$ w zbiór, który zawiera wartości dążące do zera, lecz nie zawiera zera, a więc nie jest domknięty. Ciągłość tego odwzorowania wynika z rachunku ciągłości udowodnionego poniżej.

\section{Działania na odwzorowaniach ciągłych}
\begin{twierdzenie}
\label{ma1:thm:continuous-composition}
Niech $X,Y,Z$ będą przestrzeniami metrycznymi, $A\subseteq X$, $B\subseteq Y$, $f:A\to B$, $g:B\to Z$ i $p\in A$. Jeśli $f$ jest ciągłe w $p$, a $g$ jest ciągłe w $f(p)$, to $g\circ f$ jest ciągłe w $p$. W szczególności złożenie dwóch odwzorowań ciągłych na swoich dziedzinach jest ciągłe.
\end{twierdzenie}
\begin{proof}
Dla $\varepsilon>0$ ciągłość $g$ w $f(p)$ daje $\rho>0$ takie, że dla $y\in B$ z $d_Y(y,f(p))<\rho$ mamy $d_Z(g(y),g(f(p)))<\varepsilon$. Ciągłość $f$ w $p$ daje $\delta>0$ takie, że dla $x\in A$ z $d_X(x,p)<\delta$ zachodzi $d_Y(f(x),f(p))<\rho$. Ponieważ $f(x)\in B$, otrzymujemy żądaną nierówność dla złożenia. Nie musimy wykluczać $f(x)=f(p)$.
\end{proof}

\begin{twierdzenie}
\label{ma1:thm:continuity-algebra}
Niech $X$ będzie przestrzenią metryczną, $A\subseteq X$, $K=\R$ lub $K=\C$, $p\in A$, a $f,g:A\to K$ będą ciągłe w $p$. Dla $\alpha,\beta\in K$ funkcje $\alpha f+\beta g$, $fg$ oraz $|f|$ są ciągłe w $p$. Jeżeli $g(p)\ne0$, to istnieje $\delta>0$ takie, że $g(x)\ne0$ dla wszystkich $x\in B_X(p,\delta)\cap A$, a funkcja $f/g$ na zbiorze $\{x\in A:g(x)\ne0\}$ jest ciągła w $p$.
\end{twierdzenie}
\begin{proof}
Dla dowolnego ciągu $x_n\in A$ zbieżnego do $p$ kryterium ciągowe daje $f(x_n)\to f(p)$ i $g(x_n)\to g(p)$. Rachunek granic ciągów (twierdzenie \ref{ma1:thm:limit-arithmetic}) oraz kryterium \ref{ma1:thm:continuity-sequential} dają ciągłość kombinacji liniowej, iloczynu i modułu. Jeśli $g(p)\ne0$, stosujemy ciągłość $g$ dla $\varepsilon=|g(p)|/2$; z nierówności odwrotnej otrzymujemy $|g(x)|>|g(p)|/2$ w odpowiednim otoczeniu względnym. Dla każdego ciągu w dziedzinie ilorazu zbieżnego do $p$ rachunek granic daje granicę $f(p)/g(p)$. Ponowne zastosowanie kryterium ciągowego kończy dowód.
\end{proof}

\begin{przyklad}
\label{ma1:ex:continuous-polynomials}
Funkcja $x\mapsto x$ na $K$ jest ciągła, ponieważ odległość obrazów równa się odległości argumentów. Funkcje stałe są ciągłe, ponieważ odległość ich obrazów wynosi zero. Indukcja i twierdzenie \ref{ma1:thm:continuity-algebra} dają ciągłość każdego wielomianu na $K$ i każdej funkcji wymiernej na zbiorze, na którym jej mianownik jest niezerowy.
\end{przyklad}

\begin{wniosek}
\label{ma1:cor:continuous-local-bounds}
Niech $f:A\to\R$ będzie ciągła w $p\in A$, gdzie $A$ jest podzbiorem przestrzeni metrycznej $X$. Funkcja $f$ jest ograniczona w pewnym otoczeniu względnym punktu $p$. Jeżeli $f(p)>0$, to istnieje otoczenie względne $p$, na którym $f(x)>0$. Jeżeli $f(p)<0$, to istnieje takie otoczenie, na którym $f(x)<0$.
\end{wniosek}
\begin{proof}
Z ciągłości dla $\varepsilon=1$ otrzymujemy $|f(x)|\le |f(p)|+|f(x)-f(p)|<|f(p)|+1$ w pewnym otoczeniu. Gdy $f(p)>0$, dla $\varepsilon=f(p)/2$ mamy $f(x)>f(p)-\varepsilon=f(p)/2>0$. W przypadku $f(p)<0$ stosujemy poprzedni argument do ciągłej funkcji $-f$.
\end{proof}

\section{Zwartość ciągowa i twierdzenie Weierstrassa}
\begin{definicja}
\label{ma1:def:sequentially-compact}
Niech $X$ będzie przestrzenią metryczną. Zbiór $K\subseteq X$ nazywamy ciągowo zwartym, jeżeli każdy ciąg w $K$ ma podciąg zbieżny do pewnego punktu należącego do $K$.
\end{definicja}
\begin{lemat}
\label{ma1:lem:interval-seqcompact}
Każdy przedział $[a,b]\subseteq\R$, gdzie $a,b\in\R$ i $a\le b$, jest ciągowo zwarty.
\end{lemat}
\begin{proof}
Każdy ciąg w $[a,b]$ jest ograniczony. Z twierdzenia Bolzana--Weierstrassa \ref{ma1:thm:bolzano-weierstrass} ma podciąg zbieżny do $p\in\R$. Ponieważ dla każdego wyrazu podciągu zachodzi $a\le x_{n_k}\le b$, zachowanie nieostrych nierówności w granicy daje $a\le p\le b$.
\end{proof}
\begin{twierdzenie}
\label{ma1:thm:continuous-compact-image}
Niech $X,Y$ będą przestrzeniami metrycznymi, $K\subseteq X$ ciągowo zwartym zbiorem i $f:K\to Y$ odwzorowaniem ciągłym. Wówczas $f(K)$ jest ciągowo zwarty.
\end{twierdzenie}
\begin{proof}
Dla dowolnego ciągu $(y_n)$ w $f(K)$ wybieramy $x_n\in K$ z $f(x_n)=y_n$. Ze zwartości ciągowej istnieje podciąg $x_{n_k}$ zbieżny do $p\in K$. Z ciągłości $y_{n_k}=f(x_{n_k})\to f(p)\in f(K)$. Gdy $K=\varnothing$, warunek jest spełniony w sposób pusty.
\end{proof}

\begin{twierdzenie}[Twierdzenie Weierstrassa o osiąganiu kresów]
\label{ma1:thm:weierstrass-extrema}
Niech $X$ będzie przestrzenią metryczną, $K\subseteq X$ niepustym ciągowo zwartym zbiorem i $f:K\to\R$ funkcją ciągłą. Wówczas $f$ jest ograniczona i istnieją $p,q\in K$ takie, że
\[
 f(p)=\inf f(K),\qquad f(q)=\sup f(K).
\]
\end{twierdzenie}
\begin{proof}
Gdyby $f$ nie była ograniczona, dla każdego $n\ge1$ istniałby $x_n\in K$ z $|f(x_n)|>n$. Ze zwartości ciągowej wybieramy podciąg $x_{n_k}\to p_0\in K$. Z ciągłości $f(x_{n_k})\to f(p_0)$, więc ciąg obrazów jest ograniczony. Jest to sprzeczne z $|f(x_{n_k})|>n_k\to+\infty$.

Niepusty ograniczony zbiór $f(K)$ ma kres górny $s$ i dolny $m$. Dla każdego $n\ge1$ z kryterium kresu wybieramy $u_n\in K$ z $s-1/n<f(u_n)\le s$. Zatem $f(u_n)\to s$. Wybieramy podciąg $u_{n_k}\to q\in K$. Z ciągłości $f(u_{n_k})\to f(q)$, a z wcześniejszej zbieżności także $f(u_{n_k})\to s$. Jednoznaczność granicy daje $f(q)=s$. Stosując ten sam argument do funkcji $-f$, otrzymujemy $p\in K$ z $f(p)=m$.
\end{proof}
\begin{wniosek}
\label{ma1:cor:interval-extrema}
Każda ciągła funkcja $f:[a,b]\to\R$, gdzie $a\le b$, osiąga minimum i maksimum na $[a,b]$.
\end{wniosek}
\begin{proof}
Stosujemy lemat \ref{ma1:lem:interval-seqcompact} i twierdzenie \ref{ma1:thm:weierstrass-extrema}.
\end{proof}
\begin{przyklad}
\label{ma1:ex:extrema-assumptions}
Na $(0,1)$ funkcja $f(x)=x$ jest ciągła i ograniczona, lecz nie osiąga żadnego z kresów $0,1$. Funkcja $g(x)=1/x$ na $(0,1)$ jest ciągła, lecz nieograniczona. Na domkniętym, ale nieograniczonym przedziale $[0,+\infty)$ funkcja $f(x)=x$ jest ciągła i nieograniczona. Sama ciągłość ani sama domkniętość dziedziny nie wystarcza.

Z kolei funkcja $h:[0,1]\to\R$ dana przez $h(x)=x$ dla $x<1$ i $h(1)=0$ nie osiąga kresu górnego $1$. Jest nieciągła w $1$, ponieważ $h(1-1/n)\to1$ dla $n\ge2$, lecz $h(1)=0$. Przykład pokazuje rolę ciągłości także na przedziale domkniętym i ograniczonym.
\end{przyklad}

\section{Twierdzenie o wartości pośredniej}
\begin{twierdzenie}[Własność Darboux funkcji ciągłych]
\label{ma1:thm:intermediate-value}
Niech $a,b\in\R$, $a\le b$, i niech $f:[a,b]\to\R$ będzie ciągła. Dla każdej liczby $t$ leżącej między $f(a)$ i $f(b)$, z końcami włącznie, istnieje $c\in[a,b]$ takie, że $f(c)=t$.
\end{twierdzenie}
\begin{proof}
Jeśli $t=f(a)$ lub $t=f(b)$, wybieramy odpowiedni koniec przedziału. W pozostałych przypadkach po zastąpieniu $f$ przez $-f$ i $t$ przez $-t$, jeśli to konieczne, możemy założyć $f(a)<t<f(b)$. Zbiór
$E=\{x\in[a,b]:f(x)<t\}$ jest niepusty, ponieważ $a\in E$, i ograniczony z góry przez $b$. Niech $c=\sup E\in[a,b]$.

Dla każdego $n\ge1$ wybieramy $x_n\in E$ z $c-1/n<x_n\le c$. Mamy $x_n\to c$, a z ciągłości $f(x_n)\to f(c)$. Ponieważ $f(x_n)<t$, zachodzi $f(c)\le t$. W szczególności $c\ne b$, gdyż $f(b)>t$.

Gdyby $f(c)<t$, ciągłość dla $\varepsilon=(t-f(c))/2$ dawałaby $\delta>0$ takie, że $f(x)<t$ dla $x\in[a,b]$ spełniających $|x-c|<\delta$. Liczba $x=c+\min(\delta,b-c)/2$ należy do $[a,b]$, jest większa od $c$ i spełnia $|x-c|<\delta$. Wówczas $x\in E$, co przeczy temu, że $c$ jest ograniczeniem górnym $E$. Stąd $f(c)=t$.
\end{proof}
\begin{wniosek}
\label{ma1:cor:continuous-interval-image}
Niech $I\subseteq\R$ będzie niepustym przedziałem i $f:I\to\R$ funkcją ciągłą. Wówczas $f(I)$ jest przedziałem, być może jednoelementowym. Jeżeli $I=[a,b]$, to $f(I)=[\min f(I),\max f(I)]$.
\end{wniosek}
\begin{proof}
Dla $u,v\in f(I)$ wybieramy $x,y\in I$ z $f(x)=u$, $f(y)=v$. Przedział domknięty o końcach $x,y$ zawiera się w $I$. Twierdzenie \ref{ma1:thm:intermediate-value} daje wszystkie wartości pomiędzy $u$ i $v$ w obrazie tego przedziału, a więc w $f(I)$. To jest własność charakteryzująca przedziały. Dla $I=[a,b]$ oba kresy należą do obrazu na mocy wniosku \ref{ma1:cor:interval-extrema}; żadna wartość obrazu nie leży poza nimi.
\end{proof}
\begin{przyklad}
\label{ma1:ex:ivt-root}
Wielomian $P(x)=x^3+x-1$ jest ciągły, a $P(0)=-1$ i $P(1)=1$. Twierdzenie o wartości pośredniej daje $c\in(0,1)$ z $P(c)=0$. Pierwiastek jest jedyny w $\R$: dla $x<y$ zachodzi
$P(y)-P(x)=(y-x)(y^2+xy+x^2+1)>0$, ponieważ
$y^2+xy+x^2=(y+x/2)^2+3x^2/4\ge0$.
Twierdzenie o wartości pośredniej zapewnia istnienie, natomiast powyższa nierówność zapewnia jednoznaczność.
\end{przyklad}
\begin{center}
\begin{tikzpicture}[x=6cm,y=1.4cm,>=Stealth]
 \draw[->] (-.10,0)--(1.12,0) node[right] {$x$};
 \draw[->] (0,-1.25)--(0,1.35) node[above] {$y$};
 \draw[blue!65!black,thick,domain=0:1,samples=80] plot (\x,{\x*\x*\x+\x-1});
 \fill (0,-1) circle (1.6pt) node[left] {$-1$};
 \fill (1,1) circle (1.6pt) node[right] {$1$};
 \draw[dashed,gray] (1,0) node[below] {$1$}--(1,1);
 \fill (.6823278,0) circle (2pt) node[below right] {$c$};
\end{tikzpicture}

\small Rysunek: ciągły wielomian przyjmuje wartość zero między argumentami $0$ i $1$.
\end{center}

\section{Ciągłość odwzorowania odwrotnego}
\begin{definicja}
\label{ma1:def:homeomorphism}
Niech $X,Y$ będą przestrzeniami metrycznymi. Bijekcję $f:X\to Y$ nazywamy homeomorfizmem, jeżeli zarówno $f$, jak i odwzorowanie odwrotne $f^{-1}:Y\to X$ są ciągłe.
\end{definicja}
Ciągła bijekcja nie zawsze jest homeomorfizmem. Niech $X=\R$ z metryką dyskretną $d(x,y)=1$ dla $x\ne y$, $d(x,x)=0$, a $Y=\R$ ze zwykłą metryką. Identyczność $f:X\to Y$ jest ciągła, gdyż każda kula w $X$ o promieniu $1/2$ jest singletonem. Odwrotność nie jest ciągła w żadnym punkcie $p$: ciąg $p+1/n$ dąży do $p$ w $Y$, lecz odległość tych punktów od $p$ w $X$ wynosi stale $1$.

\begin{twierdzenie}
\label{ma1:thm:compact-continuous-inverse}
Niech $X,Y$ będą przestrzeniami metrycznymi, przy czym $X$ jest ciągowo zwarta. Niech $f:X\to Y$ będzie ciągłą bijekcją. Wówczas $f^{-1}:Y\to X$ jest ciągłe.
\end{twierdzenie}
\begin{proof}
Niech $y\in Y$, $x=f^{-1}(y)$ i niech $y_n\to y$ w $Y$. Oznaczmy $x_n=f^{-1}(y_n)$. Załóżmy, że $x_n$ nie dąży do $x$. Wtedy istnieje $\varepsilon_0>0$ oraz podciąg $(x_{n_k})$ z $d_X(x_{n_k},x)\ge\varepsilon_0$ dla wszystkich $k$: wybieramy kolejno coraz większe indeksy, dla których nierówność zachodzi. Ze zwartości ciągowej ten podciąg ma dalszy podciąg zbieżny do $z\in X$. Oznaczmy wyrazy dalszego podciągu przez $u_j$. Z nierówności trójkąta mamy $\varepsilon_0\le d_X(u_j,x)\le d_X(u_j,z)+d_X(z,x)$. Ponieważ $d_X(u_j,z)\to0$, otrzymujemy $d_X(z,x)\ge\varepsilon_0$, więc $z\ne x$. Z ciągłości $f$ obrazy dalszego podciągu dążą do $f(z)$; jako podciąg $(y_n)$ dążą jednocześnie do $y$. Jednoznaczność granicy i iniektywność dają $f(z)=y=f(x)$, zatem $z=x$, sprzecznie z wcześniejszą nierównością. Stąd $x_n\to x$, a kryterium ciągowe dowodzi ciągłości $f^{-1}$.
\end{proof}
\begin{przyklad}
\label{ma1:ex:square-inverse-continuous}
Funkcja $f:[0,1]\to[0,1]$, $f(x)=x^2$, jest ciągła i ściśle rosnąca: dla $0\le x<y\le1$ mamy $y^2-x^2=(y-x)(y+x)>0$. Z twierdzenia o wartości pośredniej jej obraz zawiera $[0,1]$, a z nierówności $0\le x^2\le1$ jest równy temu przedziałowi. Jest więc bijekcją. Zwartość ciągowa $[0,1]$ i twierdzenie \ref{ma1:thm:compact-continuous-inverse} dają ciągłość odwrotności $y\mapsto\sqrt y$ na $[0,1]$, bez korzystania z rachunku różniczkowego.
\end{przyklad}

Ciągłość w każdym punkcie pozwala wybierać promień $\delta$ zależny zarówno od $\varepsilon$, jak i od punktu. W następnym rozdziale rozważymy sytuację, w której jeden promień działa jednocześnie dla wszystkich punktów dziedziny.
